\documentclass[12pt]{article}
\usepackage[utf8]{inputenc}
\usepackage[T1]{fontenc}
\usepackage{hyperref}
\title{Voro Language Resources for K-12 Curriculum Development}
\author{Compiled by Claude Sonnet 5 (claude-sonnet-5) for Andr\'as Kornai}
\date{}
\begin{document}
\maketitle

\section{Introduction}
Voro (ISO 639-3 vor; also Bena, Buna, Ebina, Woro, Yungur) is an Adamawa language of Adamawa State, Nigeria. It should not be confused with Vono (kch) or Vori (sde), which are Kainji languages of Kaduna State. Very little exists for Voro in search results, and no school materials were found.

\section{Primary Educational Resources}
None found.

\section{Linguistic and Reference Materials}
\begin{itemize}
\item Wikipedia, Voro language (Adamawa): \url{https://en.wikipedia.org/wiki/Voro_language_(Adamawa)}
\item Wikipedia, Bena--Mboi languages: \url{https://en.wikipedia.org/wiki/Bena%E2%80%93Mboi_languages}
\end{itemize}

\section{Religious and Community Resources}
Joshua Project reports Bible translation status as ``Translation Started'':
\url{https://joshuaproject.net/people_groups/20771}. Its Yungur Binna page: \url{https://joshuaproject.net/people_groups/16071}.

\section{Digital and Audio Resources}
Global Recordings Network lists Voro (Nigeria) audio: \url{https://globalrecordings.net/en/language/vor}

\section{Curriculum Development Context}
Documentation is thin. A curriculum would require new community-based work.
\end{document}